\section{Sonuç}
\indent\indent Tanzimat, Osmanlı müesseselerin modernleşmesinde kritik rol oynamıştır. Bu modernleşme çabalarını ceza hukuku alanında da görmekteyiz. Tanzimat Fermanı'nı takip eden 20 yıllık süreçte üç ceza 
kanunnamesi kaleme alınmıştır. Bunlardan ilki 1840 Ceza Kanunname-î Humayun, oldukça dar kapsamlı bir ceza kanunu denemesidir. Kapsamın darlığı, 1851 Kanunname-î Cedîd üzerinden giderilmeye çalışılsa da,
tekmilli bir ceza kanunnamesi olmaktan hayli uzaktı. 1858 Ceza Kanunname-î Humayun, 1810 Fransız ceza kanunu örnek alarak Osmanlı teb`asındaki herkesin vatandaşlık haklarını güvence altına almayı amaçlamıştır.

Tanzimat öncesinde, Osmanlı kadısı kendisine rehberlik edecek kazuistik bir kodekse sahip değildir. Elinde kendisine yardımcı olacak şer`î ve örfî bir takım kılavuzlar bulunsa da, bunlar sistemik bir kanun bütünü 
olmaktan çok uzaktı. Herhangi bir davada kararın ne olacağı tamamiyle kadıya bağlıydı. Kadı, kararını kimseye açıklamak zorunda değildi ve kararın temyizi neredeyse olanaksızdı. Bir hırsızın parmağının kesilmesi 
ya da yıllarca kürek cezasına çarptırılması aynı derecede olasıydı.

Fransız Devrimi'nin teb`a yerine vatandaşı koyması ile vatandaşlık hakları doğmuş, bu hakları güvence altına alan 1791 Fransız Ceza Kanunu ortaya çıkmıştır. Bu ceza kanunu, 18.yüzyıl Aydınlanma düşünürlerinin etkisiyle
bireyin hür iradesini dayanarak, suçun ve cezanın tanımının bila-istisna herkes için aynı olması temeline oturmuştur. Başka bir deyişle, suçun ve cezanın tanımı hakimin yorumunu gerektirmeyecek şekilde sarih olmalıydı.
Kanunun bütün vatandaşlara eşit şekilde uygulanması ancak bu şekilde sağlanabilirdi.

Tanzimat Fermanı, klasik ya da altın çağ olan Kanuni dönemine dönüşün artık mümkün olmadığını kabullenişidir. İlmiye sınıfının 150 senedir çözüm olarak ortaya sürdüğü kanun-ı kadim mefhumuna karşılık kalemiye sınıfının
çözüm önerisi bireyin hür iradesinin merkeze alındığı Batı tarzı müesseseler idi. Tanzimat doğal olarak bir ilmiye-kalemiye çatışmasını da beraberinde getirdi. Kalemiye sınıfı meydana getirdiği müesseseler ile ilmiye sınıfının 
siyaset alanını kısıtlmaya başladı. Özellikle 1858 Ceza Kanunnamesi şer`i hukuğun sorumluluğunda olan şahıs aleyhinde işlenen suçları da kapsamış ve şer`i hükümlerin dışında hükümler tanımlamıştır. Böylece etkileri belki de 
günümüze kadar devam edecek olan muhafazakar ilmiye ile seküler kalemiyenin devlet yöntemindeki mücadelesi başlıyordu.

1858 Ceza Kanunnamesi ile ceza hukuku hemen hemen şer`i hukuktan tamamen ayrılmıştır. 1850 Ticaret ve 1858 Arazi kanunnameriyle birlikte, şer`i hukuk ağırlıklı medeni ve veraset hukuku ihitva eden bir alana sıkışmıştır. 1858 Ceza 
Kanunnamesi'ni hazırlayan Meclis-i Âli-i Tanzîmat komisyonun başkanı Ahmet Cevdet Paşa'nın hazırladığı \textit{Mecelle-i Ahkaâm-ı Adliyye}(1876) ile birlikte, medeni hukuk da şer`i hukuktan ayrılacaktır.
